\section{Уравнения с частными производными первого порядка.\\
Полные интегралы. Огибающие и построение новых решений.}
\label{sec:Chapter4} \index{Chapter4}

\subsection{Полный интеграл}
Рассмотрим нелинейное дифференциальное уравнение с частными производными первого порядка:
\begin{equation} \label{eq::full}
    F(Du, u, x) = 0,
\end{equation}
где $x \in U \subset \R^n$ --- открытое множество, $u \colon U \to \R$ --- искомая функция,
$Du = (u_{x_1}, \dots, u_{x_n})$ --- её градиент по переменным $x$, 
а $F \colon \R^n \times \R \times U \to \R$ --- заданная гладкая функция.

\begin{definition}
    Функция $u(x, a) \in C^2(U \times A)$, где $a \in A \subset \R^n$ --- вектор параметров из открытого множества $A$, 
    называется \textit{полным интегралом} уравнения, если выполнены следующие условия:
    \begin{enumerate}
        \item Для любого фиксированного параметра $a \in A$ функция $u(x, a)$ 
          удовлетворяет исходному уравнению \eqref{eq::full}:
        \begin{equation}
            F(D_x u(x, a), u(x, a), x) = 0.
        \end{equation}
        \item Выполнено условие:
        \begin{equation}
            \operatorname{rk} (D_a u, D_{xa} u) = n,
        \end{equation}
        где матрица Якоби состоит из столбца $D_a u$ и матрицы смешанных производных $D_{xa} u = \left( \frac{\partial^2 u}{\partial x_i \partial a_j} \right)$.
    \end{enumerate}
\end{definition}

\subsection{Огибающие и построение новых решений}
Метод огибающих позволяет строить новые решения исходного уравнения, опираясь на известный полный интеграл (или произвольное семейство гладких решений).

\begin{proposition}
  Пусть  $\forall a\in A$ $u(x, a)$ --- удовлетворяет уравнению \eqref{eq::full}. 
  Если существует такая функция $\varphi \in C^1(U)$,
  $\varphi \colon U \to A$, 
  что для всех $x \in U$ выполняется условие:
  \begin{equation}
      D_a u(x, a) \big|_{a = \varphi(x)} = 0,
  \end{equation}
  то функция $v(x) = u(x, \varphi(x))$ также является решением 
  дифференциального уравнения \eqref{eq::full}.
  Функция $v(x)$ называется огибающей. 
\end{proposition}

\begin{proof}
    Вычислим частные производные функции $v(x)$, используя правило дифференцирования сложной функции:
    \begin{equation}
        \frac{\partial v(x)}{\partial x_i} = \frac{\partial u(x, a)}{\partial x_i} \bigg|_{a=\varphi(x)} 
        + \sum_{k=1}^n \frac{\partial u(x, a)}{\partial a_k} \bigg|_{a=\varphi(x)} \frac{\partial \varphi_k(x)}{\partial x_i}.
    \end{equation}
    В силу наложенного условия $D_a u(x, \varphi(x)) = 0$, второе слагаемое тождественно обращается в ноль. 
    Следовательно, градиент новой функции совпадает с градиентом исходного семейства при $a = \varphi(x)$:
    \begin{equation}
        D_x v(x) = D_x u(x, \varphi(x)).
    \end{equation}
    Подставляя функцию $v(x)$ и её градиент в исходное уравнение, получаем:
    \begin{equation}
        F(Dv, v, x) = F(D_x u(x, \varphi(x)), u(x, \varphi(x)), x) = 0,
    \end{equation}
    поскольку функция $u(x, a)$ является точным решением при любом фиксированном значении $a$.
\end{proof}

\subsection{Примеры полных интегралов}
\textbf{1. Уравнение Клеро}\\
Уравнение задается в виде:
\begin{equation}
    (x, Du) + f(Du) = u,
\end{equation}
где $f \colon \R^n \to \R$ --- произвольная заданная функция. 
Его полный интеграл ищется в классе линейных функций:
\begin{equation}
    u(x, a) = (a, x) + f(a), \quad a \in \R^n.
\end{equation}

\textbf{2. Уравнение Эйконала}\\
Рассмотрим уравнение, возникающее в геометрической оптике:
\begin{equation}
    |Du| = 1.
\end{equation}
Его полным интегралом служит семейство плоскостей, зависящее от $n$ параметров:
\begin{equation}
    u(x, a, b) = (a, x) + b, \quad \text{где } |a| = 1, \, b \in \R.
\end{equation}

\textbf{3. Уравнение Гамильтона-Якоби}\\
Для уравнения:
\begin{equation}
    u_t + H(D_x u) = 0,
\end{equation}
где $x \in \R^n$, $t \in \R$. Полный интеграл следующий:
\begin{equation}
    u(x, t, a, b) = (a, x) - t H(a) + b,
\end{equation}
где свободные параметры $a \in \R^n$ и $b \in \R$.

\subsection{Примеры построения огибающих}
Рассмотрим применение метода огибающих для нахождения решений уравнений 
с частными производными первого порядка на основе их известных полных интегралов.

\textbf{1. Уравнение вида $u^2(1 + |Du|^2) = 1$}\\
Для данного уравнения полный интеграл представляет собой семейство сферических 
поверхностей:
\begin{equation}
    u(x, a) = \pm \sqrt{1 - \|x - a\|^2},
\end{equation}
где $a \in \R^n$ --- вектор свободных параметров. 

Чтобы найти огибающую этого семейства, запишем условие стационарности по 
параметрам $a$, продифференцировав функцию $u(x, a)$ и приравняв полученный 
градиент к нулю:
\begin{equation}
    D_a u(x, a) = \mp \frac{a - x}{\sqrt{1 - \|x - a\|^2}} = 0.
\end{equation}
Из этого условия находим явную зависимость параметров от пространственных 
координат: $a = \varphi(x) = x$.

Подставив полученное выражение $a = x$ обратно в формулу полного интеграла, 
находим искомую огибающую функцию:
\begin{equation}
    v(x) = u(x, x) = \pm \sqrt{1 - 0} = \pm 1.
\end{equation}
Полученные константы $v(x) = 1$ и $v(x) = -1$ являются решениями 
исходного дифференциального уравнения.

\textbf{2. Уравнение Эйконала при $n = 2$}\\
Рассмотрим уравнение эйконала $|Du| = 1$ в двумерном пространстве $n=2$. 
Выберем семейство решений, где вектор параметров имеет вид 
$a = (\sin a_1, \cos a_1)$, а свободный член равен $a_2$:
\begin{equation}
    u(x, a) = x_1 \sin a_1 + x_2 \cos a_1 + a_2.
\end{equation}
Положим для простоты $a_2 = 0$. Условие стационарности по оставшемуся 
параметру $a_1$ имеет вид:
\begin{equation}
    \frac{\partial u}{\partial a_1} = x_1 \cos a_1 - x_2 \sin a_1 = 0 \implies \tan a_1 = \frac{x_1}{x_2}.
\end{equation}
Используя базовые тригонометрические тождества, выразим функции синуса и косинуса 
через координаты $x_1, x_2$: 
\begin{equation}
    \sin a_1 = \frac{x_1}{\pm \sqrt{x_1^2 + x_2^2}} = \frac{x_1}{\pm \|x\|}, \quad \cos a_1 = \frac{x_2}{\pm \sqrt{x_1^2 + x_2^2}} = \frac{x_2}{\pm \|x\|}.
\end{equation}
Подставляя эти выражения в исходное уравнение однопараметрического семейства, 
получаем функцию огибающей:
\begin{equation}
    v(x) = x_1 \left( \frac{x_1}{\pm \|x\|} \right) + x_2 \left( \frac{x_2}{\pm \|x\|} \right) = \pm \frac{x_1^2 + x_2^2}{\|x\|} = \pm \varPhi(x) = \pm \|x\|.
\end{equation}
Таким образом, огибающая представляет собой коническую поверхность 
$v(x) = \pm \|x\|$, которая является решением уравнения эйконала в любой точке, 
кроме начала координат.

\textbf{3. Уравнение Гамильтона-Якоби}\\
Рассмотрим уравнение Гамильтона-Якоби вида $u_t + H(D_x u) = 0$ с 
гамильтонианом $H(p) = \|p\|^2$. Полный интеграл при $b = 0$ имеет вид:
\begin{equation}
    u(x, t, a) = (a, x) - t \|a\|^2,
\end{equation}
где $a \in \R^n$ --- вектор параметров.

Запишем необходимое условие экстремума по вектору параметров $a$:
\begin{equation}
    D_a u(x, t, a) = x - 2t a = 0.
\end{equation}
Отсюда выражаем параметры через независимые переменные пространства и времени:
\begin{equation}
    a = \frac{x}{2t}.
\end{equation}
Подставив найденное значение $a$ в формулу исходного семейства решений, 
находим огибающую:
\begin{equation}
    v(x, t) = \left( \frac{x}{2t}, x \right) - t \left\| \frac{x}{2t} \right\|^2 = \frac{\|x\|^2}{2t} - \frac{\|x\|^2}{4t} = \frac{\|x\|^2}{4t}.
\end{equation}
Функция $v(x, t) = \frac{\|x\|^2}{4t}$ представляет собой новое решение исходного уравнения Гамильтона-Якоби.
